\src{193}
\section*{\centering\normalfont ADDITIF}

Les difficultés signalées dans la remarque de la page 13 sont actuellement
complètement résolues, grâce à une extension des théorèmes de dualité à des
variétés quelconques (à singularités arbitraires). Pour cette théorie, qui peut
se formuler dans le cadre général de la \fquote{théorie des schémas}, nous renvoyons
aux \fquote{Eléments de Géométrie algébrique} de J. DIEUDONNÉ et A. GROTHENDIECK,
actuellement en préparation. On trouvera quelques indications dans :

\begin{quote}
GROTHENDIECK (A.). - Cohomology theory of algebraic varieties,
Congrès int. des Mathématiciens, 1958, Edinburgh (A paraître).
\end{quote}

\begin{flushright}
[Avril 1959]
\end{flushright}

\begin{center}\rule{36mm}{0.5pt}\end{center}

\par\smallskip
\begingroup\footnotesize
\endgroup
